\chapter{Options américaines et exercice anticipé}\label{ch:dv:american-options-and-early-exercise}

L'action se négocie à 100.30 et passe ex-dividende demain matin, moment où elle
perdra environ le montant de son dividende, soit 1.00. Un call profondément dans la
monnaie, de prix d'exercice 80 et échéant dans un mois, figure dans les carnets de milliers
de détenteurs. Exercé ce soir, il livre une action valant 100.30 assortie de son dividende~: 20.30.
Conservé jusqu'à la date de détachement, il reste un call sur une action valant 99.30, et même avec
son mois d'optionalité restante, il vaut environ 19.56. La différence, 0.74
par action, revient à ceux qui sont de l'autre côté des calls
non exercés~: les vendeurs, assignés seulement par les détenteurs qui exercent.
Chaque trimestre, avant chaque dividende, la même décision est prise ou oubliée
sur l'ensemble du marché coté. Ce chapitre établit quand l'exercice anticipé est optimal,
à quoi ressemble la frontière entre exercer et détenir, comment dividendes et taux la déplacent,
comment valoriser les options américaines sans arbre complet, et ce que devient la décision
lors d'une véritable nuit d'échéance.

\section{Pourquoi et quand exercer par anticipation}

Une option américaine peut être exercée à tout moment avant l'échéance (livre~1,
chapitre~23). Son détenteur choisit un temps d'arrêt~; la valeur de l'option est la
plus grande espérance actualisée du gain sur l'ensemble de ces choix, le problème d'arrêt
optimal du livre~4, chapitre~10~:
\[
V_0=\sup_{\tau\le T}\E^{\mathbb Q}\bigl[e^{-r\tau}g(S_\tau)\bigr].
\]
Sur l'arbre du chapitre~2, cette borne supérieure correspond à l'induction rétrograde qui retient, à
chaque nœud, la plus grande valeur entre exercer et détenir.

\begin{definition}[Prime d'exercice anticipé]\label{def:dv:american-options-and-early-exercise:premium}
La \emph{prime d'exercice anticipé}\index{prime d'exercice anticipé} d'une option américaine
est l'écart entre sa valeur et celle de l'option européenne de même
prix d'exercice et même échéance. Elle n'est jamais négative, et elle est nulle lorsque
l'exercice anticipé n'est jamais optimal.
\end{definition}

\begin{proposition}[Quand l'exercice peut être optimal]\label{prop:dv:american-options-and-early-exercise:when}
Avec $r\ge0$~: un call américain sur une action qui ne verse aucun dividende avant l'échéance n'est jamais
exercé par anticipation, et sa prime est nulle~; un call sur une action qui verse un dividende peut être
exercé juste avant une date de détachement~; un put américain peut être exercé à tout moment
lorsque les taux sont positifs, et sa prime est strictement positive.
\end{proposition}
\begin{proof}
Call~: $C\ge S-Ke^{-r(T-t)}\ge S-K$, donc le call vaut au moins sa valeur d'exercice
immédiate, tant qu'il n'est pas exercé (chapitre~2). Entre deux dates de détachement, l'action ne verse aucun dividende, si bien que le raisonnement tient
jusqu'au moment qui précède chaque date de détachement, où l'action va perdre $D$ et la borne
devient $S-D-Ke^{-r(T-t)}$. Put~: profondément dans la monnaie, le put vaut au plus $Ke^{-r(T-t)}-S$
en tant qu'option européenne, soit moins que $K-S$ lorsque $r>0$~; exercer rapporte l'intérêt sur $K$.
\end{proof}

Le call et le put ne sacrifient pas les mêmes avantages en étant exercés~: le call renonce à
l'intérêt sur le prix d'exercice qu'il paie par anticipation et à l'assurance de son optionalité, et ne gagne que
les dividendes~; le put renonce à son optionalité et gagne l'intérêt sur le prix d'exercice
qu'il reçoit par anticipation. La \Cref{fig:dv:american-options-and-early-exercise:premium} montre la
prime du put en fonction du taux~: nulle à taux nul, 0.52 à 5~\%, 1.06 à 10~\% pour
un put à la monnaie d'un an, avec une volatilité de 20.

\begin{omfigure}
\begin{tikzpicture}
\begin{axis}[width=10.5cm,height=4.8cm,
  xlabel={taux d'intérêt (\%)},ylabel={prime d'exercice anticipé},
  tick label style={font=\small},label style={font=\small},
  xmin=0,xmax=10,ymin=0,ymax=1.15,grid=major,grid style={gray!25},
  table/col sep=comma]
\addplot[omThm,very thick,mark=*,mark size=1.5pt] table[x=rate,y=premium]{figdata/derivatives/06-american-options-and-early-exercise/premium.csv};
\end{axis}
\end{tikzpicture}

\omcaption{\omterm{def:dv:american-options-and-early-exercise:premium}{Prime d'exercice anticipé} d'un put américain d'un an à la monnaie
(cours au comptant 100, volatilité 20~\%, sans dividende) selon le taux d'intérêt~: la valeur de recevoir le
prix d'exercice par anticipation. À taux nul, il n'y a rien à gagner à exercer. Données~: le
tutoriel.}\label{fig:dv:american-options-and-early-exercise:premium}
\end{omfigure}

\section{La frontière d'exercice}

\begin{definition}[Frontière d'exercice]\label{def:dv:american-options-and-early-exercise:boundary}
La \emph{frontière d'exercice}\index{frontière d'exercice} d'une option américaine est la courbe
$t\mapsto S^*(t)$ qui sépare les cours au comptant où exercer est optimal de ceux où
conserver l'option l'est~: pour un put, l'exercice est optimal lorsque $S_t\le S^*(t)$. Sur la
frontière, l'option égale sa valeur d'exercice et, dans le \omterm{def:dv:black-scholes-three-ways:model}{modèle de Black--Scholes}, la rejoint avec la même
pente (recollement C1 (smooth pasting), livre~4, chapitre~10).
\end{definition}

Pour un put sans dividende, la frontière remonte vers le prix d'exercice à mesure que l'échéance approche
(lorsque $r>0$), et lorsque le temps jusqu'à l'échéance s'allonge, elle descend vers la frontière du
put perpétuel. Cette limite admet une forme fermée.

\begin{proposition}[Put américain perpétuel]\label{prop:dv:american-options-and-early-exercise:perpetual}
Dans le \omterm{def:dv:black-scholes-three-ways:model}{modèle de Black--Scholes} avec $r>0$ et sans dividende, un put qui n'expire jamais vaut
$(K-S^*)(S/S^*)^{-\gamma}$ pour $S>S^*$ et $K-S$ en dessous, avec $\gamma=2r/\sigma^2$ et
$S^*=K\gamma/(1+\gamma)$.
\end{proposition}
\begin{proof}
Sans dépendance en temps, l'équation de valorisation s'écrit $\tfrac12\sigma^2S^2V^{\prime\prime}+rSV^{\prime}-rV=0$,
résolue par $S^{-\gamma}$ et $S$~; un put qui s'annule à l'infini ne conserve que $AS^{-\gamma}$.
Le raccordement des valeurs $AS^{*-\gamma}=K-S^*$ et le recollement C1 $-\gamma AS^{*-\gamma-1}=-1$ donnent
$S^*=K\gamma/(1+\gamma)$.
\end{proof}

Avec $r=5\%$ et $\sigma=20\%$, $\gamma=2.5$ et la frontière perpétuelle vaut 71.43~; la
frontière d'un put d'un an vaut 81.11 aujourd'hui, et celle d'un put de trois ans vaut environ 76.2
(voir la \cref{fig:dv:american-options-and-early-exercise:boundary}).

\begin{omfigure}
\begin{tikzpicture}
\begin{axis}[width=10.5cm,height=5cm,
  xlabel={temps jusqu'à l'échéance (années)},ylabel={cours au comptant},
  tick label style={font=\small},label style={font=\small},
  xmin=0,xmax=3,ymin=65,ymax=104,grid=major,grid style={gray!25},
  legend columns=3,legend style={font=\footnotesize,at={(0.5,-0.3)},anchor=north,draw=none,fill=none,/tikz/every even column/.append style={column sep=8pt}},
  table/col sep=comma]
\addplot[name path=bd,omThm,very thick] table[x=tau,y=boundary]{figdata/derivatives/06-american-options-and-early-exercise/boundary.csv};
\path[name path=axisbottom] (axis cs:0,65) -- (axis cs:3,65);
\addplot[omThm!15] fill between[of=bd and axisbottom];
\addplot[omDef,thick,dashed] table[x=tau,y=perpetual]{figdata/derivatives/06-american-options-and-early-exercise/boundary.csv};
\addplot[gray,thick,dotted] table[x=tau,y=strike]{figdata/derivatives/06-american-options-and-early-exercise/boundary.csv};
\node[font=\footnotesize,text=omThm] at (axis cs:2.2,68) {exercice};
\node[font=\footnotesize] at (axis cs:2.2,92) {détention};
\legend{frontière d'exercice,,limite perpétuelle 71.43,prix d'exercice}
\end{axis}
\end{tikzpicture}

\omcaption{\omterm{def:dv:american-options-and-early-exercise:boundary}{Frontière d'exercice} d'un put américain (prix d'exercice 100, $r=5\%$, $\sigma=20\%$) en
fonction du temps restant avant l'échéance, calculée sur un arbre de 3\,000 pas~: sous la courbe, le
détenteur exerce. La frontière remonte vers le prix d'exercice à mesure que l'échéance approche et
descend vers les 71.43 du put perpétuel lorsque l'horizon s'allonge. Les pas sont ceux de la grille
de l'arbre. Données~: le tutoriel.}\label{fig:dv:american-options-and-early-exercise:boundary}
\end{omfigure}

\section{Dividendes et taux comme déclencheurs}

Pour un call, la seule raison d'exercer par anticipation est de capturer un dividende, et le seul
moment est juste avant une date de détachement. Le détenteur compare deux valeurs la veille au soir.

\begin{method}[La décision d'exercice d'un call avant une date de détachement]\label{met:dv:american-options-and-early-exercise:decision}
Avec un spot $S$, un prix d'exercice $K$, un dividende $D$ et un temps $T-t$ restant après la date de détachement~:
\begin{enumerate}
\item valeur d'exercice~: $S-K$ (l'action, avec son dividende)~;
\item valeur de détention~: au moins le call européen sur l'action ex-dividende,
$C(S-D,K,T-t)$~;
\item on exerce si la première dépasse la seconde. Comme $C(S-D,K,T-t)\ge S-D-Ke^{-r(T-t)}$,
l'exercice exige $D>K\bigl(1-e^{-r(T-t)}\bigr)$~: le dividende doit dépasser l'intérêt
sur le prix d'exercice pendant la durée de vie restante. Pour un call profondément dans la monnaie,
dont la valeur temps ne dépasse guère cet intérêt, le seuil lui est proche.
\end{enumerate}
\end{method}

\begin{example}[La veille de la date de détachement]\label{ex:dv:american-options-and-early-exercise:night}
Spot 100.30, prix d'exercice 80, 29 jours restants, $r=4\%$, $\sigma=25\%$, dividende 1.00. L'exercice
vaut 20.30, la détention 19.56. L'intérêt sur le prix d'exercice pendant 29 jours est de 0.254, et le
dividende au-delà duquel l'exercice est gagnant est de 0.255~: la composante de put dans la valeur temps
du call vaut un dixième de centime. Avec un dividende de 1.00, le détenteur qui n'exerce pas abandonne
0.74 par action (\cref{fig:dv:american-options-and-early-exercise:decision}).
\end{example}

\begin{omfigure}
\begin{tikzpicture}
\begin{axis}[width=10.5cm,height=4.8cm,
  xlabel={dividende versé demain},ylabel={valeur du call 80 ce soir},
  tick label style={font=\small},label style={font=\small},
  xmin=0,xmax=1,ymin=19.2,ymax=20.6,grid=major,grid style={gray!25},
  legend columns=2,legend style={font=\footnotesize,at={(0.5,-0.3)},anchor=north,draw=none,fill=none,/tikz/every even column/.append style={column sep=8pt}},
  table/col sep=comma]
\addplot[omThm,very thick] table[x=div,y=exercise]{figdata/derivatives/06-american-options-and-early-exercise/decision.csv};
\addplot[omDef,very thick] table[x=div,y=hold]{figdata/derivatives/06-american-options-and-early-exercise/decision.csv};
\draw[gray,dashed] (axis cs:0.255,19.2) -- (axis cs:0.255,20.6);
\node[font=\footnotesize,anchor=west] at (axis cs:0.27,19.35) {seuil 0.255};
\legend{exercice ce soir,détention jusqu'à la date de détachement}
\end{axis}
\end{tikzpicture}

\omcaption{Le call de l'\cref{ex:dv:american-options-and-early-exercise:night} la veille de la date de
détachement~: sa valeur d'exercice ne dépend pas du dividende, sa valeur conservée jusqu'à la date de
détachement baisse d'un pour un avec lui. Au-delà d'un dividende de 0.255, le détenteur doit exercer.
Données~: le tutoriel.}\label{fig:dv:american-options-and-early-exercise:decision}
\end{omfigure}

Les puts évoluent en sens inverse face aux dividendes~: un dividende abaisse l'action après la date de
détachement et augmente donc le put, ce qui rend la détention d'un put à travers une date de détachement plus attrayante et
l'exercice anticipé juste avant moins attrayant. Des taux élevés et une option profondément dans la monnaie rendent l'exercice anticipé des
puts plus probable~; une forte volatilité le rend moins probable, car elle accroît la valeur
de l'attente, et il en va de même d'un rendement de dividende ou de frais d'emprunt (chapitre 5), qui
abaisse le forward que le put attend. Les frais d'emprunt jouent en sens inverse pour les calls~: un détenteur de call
qui l'exerce possède une action pouvant être prêtée moyennant ces frais, et sur une action difficile à emprunter, cela peut
justifier d'exercer un call par anticipation sans aucun dividende en vue.

\section{Approximations et bornes}

Un arbre price une option américaine en quelques millisecondes. Un teneur de marché qui recalcule des
dizaines de milliers de prix à chaque tick a besoin de quelque chose de plus rapide, et deux
approximations en forme fermée sont courantes.

\begin{definition}[Approximation quadratique]\label{def:dv:american-options-and-early-exercise:baw}
L'\emph{approximation quadratique}\index{approximation quadratique} de Barone-Adesi et Whaley
écrit la valeur américaine comme la valeur européenne plus une prime de la forme $A(S/S^*)^{q}$,
où l'exposant $q$ est solution de l'équation du second degré qu'impose la partie indépendante du temps de l'équation
de pricing, et où le spot critique $S^*$ se détermine numériquement par égalité des valeurs
et recollement C1 (smooth pasting) en $S^*$. Au-delà de $S^*$, l'option vaut sa valeur d'exercice.
\end{definition}

L'approximation est exacte pour l'option perpétuelle et pour les échéances très courtes, et bonne
entre les deux (voir le tableau ci-dessous). L'approximation de Bjerksund--Stensland repose sur une idée
différente~: elle price exactement la stratégie qui exerce à un déclencheur plat, laquelle est réalisable
mais non optimale, et fournit donc une borne inférieure proche de la valeur vraie.

\begin{center}
\footnotesize
\begin{tabular}{@{}rrrrrrrr@{}}
\toprule
spot & échéance & taux & rendement & européen & arbre (2\,001) & quadratique & erreur \\
\midrule
100 & 1 & 5~\% & 0 & 5.5735 & 6.0902 & 6.0976 & $+0.0074$ \\
90 & 0.5 & 8~\% & 0 & 11.2671 & 12.1792 & 12.1151 & $-0.0641$ \\
110 & 0.25 & 3~\% & 0 & 1.5070 & 1.5201 & 1.5229 & $+0.0028$ \\
100 & 1 & 5~\% & 4~\% & 9.0396 & 9.2168 & 9.2444 & $+0.0277$ \\
\bottomrule
\end{tabular}

\smallskip
Puts américains de prix d'exercice 100 (volatilité 20~\%, 30~\%, 25~\%, 25~\% selon la ligne). L'\omterm{def:dv:american-options-and-early-exercise:baw}{approximation
quadratique} reste à moins de trois centimes, sauf profondément dans la monnaie à taux élevé.
\end{center}

\begin{definition}[Exercice bermudéen]\label{def:dv:american-options-and-early-exercise:bermudan}
Une option à \emph{exercice bermudéen}\index{exercice bermudéen} ne peut être exercée qu'à un ensemble
fini de dates avant l'échéance. Sa valeur se situe entre les valeurs européenne et américaine et croît
avec l'ensemble des dates d'exercice.
\end{definition}

L'\omterm{def:dv:american-options-and-early-exercise:bermudan}{exercice bermudéen} est la règle des obligations remboursables par anticipation (callable) et des swaptions
bermudéennes du livre~6, chapitre~9, et c'est ce que chaque méthode numérique price en pratique~: un
arbre à $n$ pas est un bermudéen à $n$ dates. La valeur converge vers la valeur américaine à mesure que
les dates se multiplient, rapidement au début~: un put d'un an exerçable chaque mois vaut 6.04,
déjà à 90~\% du chemin entre la valeur européenne 5.57 et la valeur américaine 6.09
(\cref{fig:dv:american-options-and-early-exercise:bermudan}).

\begin{omfigure}
\begin{tikzpicture}
\begin{semilogxaxis}[width=10.5cm,height=4.8cm,
  xlabel={nombre de dates d'exercice (échelle log)},ylabel={valeur du put},
  tick label style={font=\small},label style={font=\small},
  xmin=1,xmax=200,ymin=5.5,ymax=6.15,grid=major,grid style={gray!25},
  legend columns=3,legend style={font=\footnotesize,at={(0.5,-0.3)},anchor=north,draw=none,fill=none,/tikz/every even column/.append style={column sep=8pt}},
  table/col sep=comma]
\addplot[omThm,very thick,mark=*,mark size=1.5pt] table[x=dates,y=bermudan]{figdata/derivatives/06-american-options-and-early-exercise/bermudan.csv};
\addplot[omDef,thick,dashed] table[x=dates,y=european]{figdata/derivatives/06-american-options-and-early-exercise/bermudan.csv};
\addplot[omProp,thick,dotted] coordinates {(1,6.0902) (200,6.0902)};
\legend{Bermudéen,Européen 5.5735,Américain 6.0902}
\end{semilogxaxis}
\end{tikzpicture}

\omcaption{Un put d'un an à la monnaie ($r=5\%$, $\sigma=20\%$) exerçable sur 1 à 200
dates également espacées, la dernière à l'échéance. Une seule date donne l'option européenne~; l'exercice
mensuel capture 90~\% de la \omterm{def:dv:american-options-and-early-exercise:premium}{prime d'exercice anticipé}. Données~: le
tutoriel.}\label{fig:dv:american-options-and-early-exercise:bermudan}
\end{omfigure}

\section{La décision pratique d'exercice}

La théorie dit quand~; le marché décide si. L'exercice est une instruction envoyée par le
détenteur, par l'intermédiaire du courtier, à la chambre de compensation avant une heure limite~; seules les options au moins
d'un cent dans la monnaie à l'échéance sont exercées automatiquement (livre~1, chapitre~23),
et rien n'est automatique avant l'échéance. Un détenteur qui oublie d'exercer un call la veille
d'une date de détachement perd l'écart calculé plus haut, et l'assignation
aléatoire de la chambre de compensation la transmet aux vendeurs non assignés.

\begin{dated}{2026-09}{Calls non exercés avant les dates de détachement}\label{dat:dv:american-options-and-early-exercise:pool}
Une étude portant sur des calls sur actions américaines versant des dividendes trimestriels, de janvier 1996 à
avril 2006, a montré que plus de la moitié des positions longues qui auraient dû être
exercées la veille d'une date de détachement ne l'ont pas été, ce qui a coûté aux détenteurs plus de
491~millions de dollars américains sur la période~; les teneurs de marché en ont capté l'essentiel
grâce à une stratégie de spread sur dividende (chapitre~26). Le seuil d'exercice automatique d'un cent
à l'échéance est la règle de la chambre de compensation des options sur actions et sur indices.
\end{dated}

Trois points pratiques en découlent. Les desks évaluent leurs calls américains vendus en supposant que
les détenteurs exercent de façon optimale, et enregistrent le gain lorsqu'ils n'exercent pas~; valoriser la
position du vendeur avec un taux de défaillance supposé est un risque de modèle (chapitre~27). Les volatilités
implicites des options américaines se calculent avec un pricer américain ou après retrait de la
\omterm{def:dv:american-options-and-early-exercise:premium}{prime d'exercice anticipé}~: inverser la formule européenne au prix d'un put américain profondément dans
la monnaie donne une volatilité trop élevée. Les décisions proches de l'échéance se prennent avec
le dividende et les frais d'emprunt du jour, qui déplacent les seuils autant que le spot.

\section{Tutoriel~: frontières, approximations et une décision}\label{tut:dv:american-options-and-early-exercise:tut}

\begin{tutorial}
\textbf{Objectif.} Calculer la \omterm{def:dv:american-options-and-early-exercise:boundary}{frontière d'exercice} du put sur un arbre, comparer l'\omterm{def:dv:american-options-and-early-exercise:baw}{approximation quadratique}
à l'arbre, valoriser l'\omterm{def:dv:american-options-and-early-exercise:bermudan}{exercice bermudéen}, et résoudre la décision relative au call
avant une date de détachement. \textbf{État final~:} les quatre figures du chapitre, le tableau et les
chiffres de l'\cref{ex:dv:american-options-and-early-exercise:night}.
\begin{tutsteps}
\item \textbf{La frontière}~: à chaque pas de l'arbre, le nœud le plus élevé où l'exercice
est avantageux.
\omcode{code/firm/american/firm_american.py}{63}{84}{La frontière d'exercice du put lue sur un arbre de Cox--Ross--Rubinstein.}\label{lst:dv:american-options-and-early-exercise:boundary}
\item \textbf{La décision} et son seuil de dividende~:
\omcode{code/firm/american/firm_american.py}{106}{125}{Exercer ou conserver la position la veille d'une date de détachement, et le dividende-seuil.}\label{lst:dv:american-options-and-early-exercise:decision}
\item \textbf{Exécuter} \texttt{dv\_american.results()}, \texttt{accuracy\_table()},
\texttt{hook()} et \texttt{fig\_american.py}.
\end{tutsteps}
\textbf{Pour aller plus loin.} Ajouter un second dividende deux semaines après le premier et déterminer le
seuil de dividende avant chaque date de détachement~; calculer ensuite la frontière du put avec un
rendement de dividende de 4~\% et la comparer à celle du call.
\end{tutorial}

\section{Construction~: le pricer américain et l'assistant d'exercice}\label{bld:dv:american-options-and-early-exercise:american}

\begin{build}
\bfield{Rôle} Le desk d'options cotées de la firme miniature valorise ses options américaines
avec l'approximation de la boucle de cotation et avec l'arbre pour le risque~; à chaque veille de
date de détachement, il recense les calls, longs et courts, dont l'exercice est optimal.
\bfield{Interface} \texttt{baw(spot, strike, t, r, q, vol, right) -> (price, critical)}~;
\texttt{put\_boundary(strike, t, r, q, vol, n)}~; \texttt{bermudan\_put(\ldots, n\_ex)}~;
\texttt{exercise\_decision(spot, strike, dividend, t\_left, r, vol)}~;
\texttt{dividend\_threshold(\ldots)}~; l'arbre lui-même est \texttt{firm.binomial}.
\bfield{Règles} Le call américain sans dividende égale le call européen~; l'approximation
renvoie la valeur d'exercice au-delà de son spot critique~; la valeur de détention dans la décision est
le call européen sur l'action détachée.
\bfield{Tests d'acceptation} \texttt{code/firm/american/tests/}~: l'approximation à moins
de 0.07 d'un arbre à 2\,001 pas sur les quatre cas du tableau~; la limite perpétuelle de la
frontière~; la valeur bermudéenne croissante avec le nombre de dates, comprise entre la valeur européenne et
la valeur américaine~; le dividende seuil juste au-dessus de l'intérêt sur le strike.
\bfield{Pour aller plus loin} La borne inférieure de Bjerksund--Stensland, et une décision tenant compte des dividendes qui
valorise la détention avec le pricer américain lorsqu'un dividende supplémentaire suit.
\end{build}

\begin{omsources}
\item R.~C.~Merton, ``Theory of rational option pricing'', \emph{Bell Journal of Economics
and Management Science} 4 (1973) 141--183.
\item G.~Barone-Adesi and R.~E.~Whaley, ``Efficient analytic approximation of American
option values'', \emph{Journal of Finance} 42 (1987) 301--320.
\item P.~Bjerksund and G.~Stensland, ``Closed-form approximation of American options'',
\emph{Scandinavian Journal of Management} 9 (1993) 87--99.
\item V.~K.~Pool, H.~R.~Stoll and R.~E.~Whaley, ``Failure to exercise call options: an
anomaly and a trading game'', \emph{Journal of Financial Markets} 11 (2008).
\end{omsources}

\section{Exercices}

\begin{exercise}[$\star$]\label{exo:dv:american-options-and-early-exercise:1}
Un put américain d'un an de prix d'exercice 100 sur une action à 100 ($r=5\%$, $\sigma=20\%$) vaut
6.0902 et la valeur européenne 5.5735. Donner la \omterm{def:dv:american-options-and-early-exercise:premium}{prime d'exercice anticipé}, et sa taille relative
à la valeur européenne.
\end{exercise}

\begin{exercise}[$\star$]\label{exo:dv:american-options-and-early-exercise:2}
Donner la frontière et la valeur au spot 100 d'un put américain perpétuel de prix d'exercice 100 avec
$r=5\%$ et $\sigma=20\%$.
\end{exercise}

\begin{exercise}[$\star$]\label{exo:dv:american-options-and-early-exercise:3}
Un call de prix d'exercice 50 a 60 jours restants après le détachement de demain~; $r=5\%$. En dessous de quel dividende
son exercice anticipé n'est-il certainement pas optimal~?
\end{exercise}

\begin{exercise}[$\star\star$]\label{exo:dv:american-options-and-early-exercise:4}
Pourquoi une forte volatilité rend-elle l'exercice anticipé d'un put moins probable, et un taux
d'intérêt élevé plus probable~? Quel rôle jouent les frais d'emprunt~?
\end{exercise}

\begin{exercise}[$\star\star$]\label{exo:dv:american-options-and-early-exercise:5}
Montrer que le put bermudéen dont la seule date d'exercice est l'échéance est le put européen, et que
l'ajout de dates ne peut pas diminuer sa valeur.
\end{exercise}

\begin{exercise}[$\star\star$]\label{exo:dv:american-options-and-early-exercise:6}
D'après le tableau, dans quel cas l'\omterm{def:dv:american-options-and-early-exercise:baw}{approximation quadratique} s'écarte-t-elle le plus du prix exact, et pourquoi à cet endroit~?
\end{exercise}

\begin{exercise}[$\star\star\star$]\label{exo:dv:american-options-and-early-exercise:7}
\emph{Codage.} Avec \texttt{dividend\_threshold}, trouver le dividende seuil du call de
l'\cref{ex:dv:american-options-and-early-exercise:night} avec 90 jours restants au lieu de 29, et
le comparer à l'intérêt sur le prix d'exercice.
\end{exercise}

\begin{exercise}[$\star\star\star$]\label{exo:dv:american-options-and-early-exercise:8}
\emph{Repérer l'erreur.} «~Le put américain cote 6.09. En inversant la formule européenne
à 6.09, on obtient une volatilité de 21.4~\%, contre 20~\% pour les options européennes. Les options
américaines sont chères~: vendez-les.~»
\end{exercise}

\section{Problème~: la veille du détachement}

\begin{problem}[{Problème du week-end --- qui perçoit le dividende}]\label{pb:dv:american-options-and-early-exercise:1}
L'action s'échange à 100.30 la veille de son détachement d'un dividende
de 1.00. Le call 80 a 29 jours restants après demain, $r=4\%$, une volatilité de 25~\%, et des
positions ouvertes de 10\,000 contrats (multiplicateur 100). L'expérience suggère que 45~\% des détenteurs
n'exerceront pas ce soir.

\medskip\noindent\textbf{Partie I --- La décision.}
\begin{enumerate}
\item Donner la valeur d'exercice ce soir.
\item Donner la valeur de la détention jusqu'à la date de détachement.
\item Un détenteur doit-il exercer~? De combien par action~?
\item Donner l'intérêt sur le prix d'exercice pendant les 29 jours restants.
\item Donner le dividende seuil, et expliquer pourquoi il est si proche de l'intérêt.
\end{enumerate}

\medskip\noindent\textbf{Partie II --- Qui paie.}
\begin{enumerate}[resume]
\item Donner la perte de l'ensemble des détenteurs si personne n'a exercé.
\item Donner le transfert aux vendeurs si 45~\% des détenteurs n'exercent pas.
\item Quels vendeurs le reçoivent, et comment l'assignation aléatoire le décide-t-elle~?
\item Pourquoi un teneur de marché vendeur de ces calls ne peut-il pas compter sur ce transfert~?
\item Que détient un vendeur assigné demain matin, et que doit-il~?
\end{enumerate}

\medskip\noindent\textbf{Partie III --- Prix.}
\begin{enumerate}[resume]
\item Où doit se situer le cours acheteur du call ce soir, par rapport à 20.30~?
\item Demain, que vaut le call si l'action ouvre à 99.30~?
\item Un trader achète l'action et vend le call ce soir à 20.30. Que rapporte la transaction
si le call n'est pas exercé contre elle~?
\item Et s'il l'est~?
\item Comment dimensionner une telle transaction compte tenu de l'incertitude sur le taux de non-exercice~?
\end{enumerate}

\medskip\noindent\textbf{Partie IV --- Jugement.}
\begin{enumerate}[resume]
\item Pourquoi les détenteurs particuliers n'exercent-ils pas~?
\item Un courtier doit-il exercer automatiquement pour ses clients~?
\item Comment un desk doit-il valoriser ses calls vendus ce soir~?
\item Énoncer le \emph{résultat nommé}~: le dividende seuil et la valeur transférée aux
vendeurs par les détenteurs qui n'exercent pas.
\item En une phrase~: quand un call américain vaut-il plus qu'un call européen~?
\end{enumerate}
\end{problem}

\section{Questions d'entretien}

\begin{interviewq}[$\star$ \iqroles{trader}]\label{iq:dv:american-options-and-early-exercise:1}
Exerceriez-vous jamais un call américain par anticipation~? Un put américain~?
\end{interviewq}

\begin{interviewq}[$\star$ \iqroles{chercheur}]\label{iq:dv:american-options-and-early-exercise:2}
Pourquoi la valeur d'une option américaine est-elle un supremum sur les temps d'arrêt, et comment un arbre
la calcule-t-il~?
\end{interviewq}

\begin{interviewq}[$\star\star$ \iqroles{chercheur}]\label{iq:dv:american-options-and-early-exercise:3}
Dérivez la \omterm{def:dv:american-options-and-early-exercise:boundary}{frontière d'exercice} d'un put américain perpétuel.
\end{interviewq}

\begin{interviewq}[$\star\star$ \iqroles{trader, risque}]\label{iq:dv:american-options-and-early-exercise:4}
Vous êtes vendeur de calls profondément dans la monnaie la veille d'un gros dividende. Qu'attendez-vous,
et comment vous couvrez-vous~?
\end{interviewq}

\begin{interviewq}[$\star\star$ \iqroles{développeur}]\label{iq:dv:american-options-and-early-exercise:5}
Vous devez valoriser 50\,000 options américaines par seconde. Arbre, grille ou approximation~? Que
vérifiez-vous~?
\end{interviewq}

\begin{interviewq}[$\star\star\star$ \iqroles{chercheur, trader}]\label{iq:dv:american-options-and-early-exercise:6}
Comment calculez-vous une volatilité implicite à partir du prix d'une option américaine, et que se passe-t-il
si vous utilisez la formule européenne~?
\end{interviewq}
